\documentclass{article}
\usepackage[final]{neurips_2025}
\usepackage[utf8]{inputenc}
\usepackage[T1]{fontenc}
\usepackage{hyperref}
\usepackage{url}
\usepackage{booktabs}
\usepackage{amsfonts}
\usepackage{amsmath}
\usepackage{amssymb}
\usepackage{nicefrac}
\usepackage{microtype}
\usepackage{graphicx}
\usepackage{xcolor}
\usepackage{multirow}
\usepackage{algorithm}
\usepackage{algorithmic}

\newcommand{\ours}{EGTM}
\newcommand{\rbase}{r_0}

\title{Entropy-Guided Adaptive Token Merging for Robust and Efficient Vision Transformers}

\author{
  Anonymous Author(s)
}

\begin{document}

\maketitle

\begin{abstract}
Token merging methods like ToMe achieve significant speedups for Vision Transformers (ViTs) by gradually merging similar tokens during inference. However, their robustness under distribution shift---the conditions most relevant to real-world deployment---has never been systematically studied. We present the first comprehensive evaluation of token reduction robustness under common corruptions (ImageNet-C) and find that aggressive token merging causes \emph{disproportionate} accuracy degradation under distribution shift. On DeiT-B with merging ratio $r{=}16$, ToMe's relative accuracy drop under corruption exceeds that of the uncompressed model by 11.6 percentage points---a gap far larger than the 1.5 percentage points observed on the smaller DeiT-S with $r{=}8$. To address this, we propose Entropy-Guided Token Merging (\ours{}), a training-free method that adapts the merging ratio based on per-layer CLS-row attention entropy. When attention entropy indicates distribution shift, \ours{} reduces merging to preserve accuracy. On DeiT-B, \ours{} fully recovers the disproportionate degradation gap (reducing it from 11.6\% to $-0.2\%$) while maintaining a $1.15\times$ throughput improvement. We find that \ours{}'s primary value lies in larger models with aggressive merging ratios; on smaller models with conservative ratios, the degradation gap is modest and simpler fixed-ratio reduction suffices. We release our benchmark and code to facilitate future research on robust token reduction.
\end{abstract}

\section{Introduction}

Vision Transformers (ViTs)~\citep{dosovitskiy2021image} have become the dominant architecture for image recognition. However, their quadratic attention complexity motivates extensive work on \emph{token reduction}---methods that prune or merge tokens during inference. Token Merging (ToMe)~\citep{bolya2023token} is the most widely adopted approach, using bipartite soft matching on key cosine similarity to merge similar tokens with a fixed ratio per layer, achieving $2\times$+ throughput with minimal accuracy loss on clean data. Recent surveys argue that token reduction should transcend its traditional efficiency role, influencing robustness and alignment~\citep{kong2025token}.

Despite their widespread deployment, \textbf{token reduction methods have never been systematically evaluated under distribution shift}---the very conditions that deployed vision systems must handle. Real-world inputs regularly contain corruptions (noise, blur, weather artifacts), and we hypothesize that token merging suffers \emph{disproportionate} accuracy degradation under these conditions for three reasons: (1) input perturbations distort the cosine similarity structure used for bipartite matching~\citep{guo2023robustifying}; (2) token merging can amplify the ``token overfocusing'' phenomenon where attention concentrates on non-robust tokens under corruption; and (3) fixed merging ratios fail to account for input difficulty.

We propose \textbf{Entropy-Guided Token Merging (\ours{})}, a training-free, drop-in replacement for ToMe that adapts the merging ratio based on per-layer CLS-row attention entropy. Under clean inputs where attention is sharp (low entropy), \ours{} merges aggressively for efficiency; under corrupted inputs where attention becomes diffuse (high entropy), it reduces merging to preserve accuracy. The method requires only a lightweight calibration step on clean data and adds negligible computational overhead.

Our contributions are:
\begin{itemize}
    \item We present the \textbf{first systematic study} of token reduction robustness under distribution shift, evaluating five token reduction methods across 12 corruption types on DeiT-S and DeiT-B. We find that aggressive token merging (DeiT-B, $r{=}16$) causes an 11.6\% disproportionate degradation gap, while conservative merging (DeiT-S, $r{=}8$) shows a modest 1.5\% gap.
    \item We propose \ours{}, a \textbf{training-free, shift-aware} token merging method that uses CLS-row attention entropy to adapt the merging ratio at inference time.
    \item On DeiT-B, \ours{} \textbf{fully recovers} the disproportionate degradation (gap: $-0.2\%$) while retaining a $1.15\times$ speedup. We honestly characterize the limitations: on DeiT-S, the entropy computation overhead exceeds merging savings, and a fixed reduced ratio ($r{=}4$) achieves comparable robustness with better throughput.
\end{itemize}

Section~\ref{sec:related} reviews related work. Section~\ref{sec:method} describes \ours{}. Section~\ref{sec:experiments} presents experiments. Section~\ref{sec:discussion} discusses limitations and implications.

\section{Related Work}
\label{sec:related}

\paragraph{Token Reduction for Vision Transformers.}
Token reduction methods broadly fall into pruning and merging approaches. EViT~\citep{liang2022evit} keeps top-$k$ attentive tokens and fuses the rest; DynamicViT~\citep{rao2021dynamicvit} learns importance prediction modules; A-ViT~\citep{yin2022avit} uses adaptive halting per token. ToMe~\citep{bolya2023token} introduced training-free bipartite soft matching for token merging and remains the most widely adopted method. Token Fusion~\citep{kim2024token} bridges pruning and merging. None of these methods study robustness under distribution shift.

\paragraph{Image-Adaptive Token Merging.}
AToM~\citep{shin2025atom} proposes image-adaptive, fine-grained token merging with algorithm-architecture co-design. AdapTiV~\citep{yoo2024adaptiv} uses sign-similarity based image-adaptive merging with hardware co-design. Both adapt \emph{which} tokens to merge for better clean-data efficiency but do not study behavior under distribution shift. In contrast, \ours{} adapts \emph{how many} tokens to merge based on a shift-detection signal. The approaches are complementary.

\paragraph{Robustness of Vision Transformers.}
Zhang et al.~\citep{zhang2022delving} showed ViTs generalize better than CNNs under distribution shift. Guo et al.~\citep{guo2023robustifying} identified ``token overfocusing'' under corruption and proposed attention diversification during training. SATA~\citep{nikzad2025sata} enhances ViT robustness through spatial autocorrelation analysis of tokens before the FFN block, without retraining. Unlike SATA, which modifies FFN computation, we address degradation of token \emph{merging} decisions. The two approaches are complementary.

\paragraph{Entropy in Vision Transformers.}
Zhai et al.~\citep{zhai2023stabilizing} established attention entropy as a diagnostic of transformer behavior, showing that entropy collapse correlates with training instability. Lin et al.~\citep{lin2026entropy} use transfer entropy for layer-wise condensation during training. We use CLS-row attention entropy at \emph{inference time} for a different purpose: detecting distribution shift to adapt token merging.

\paragraph{Test-Time Adaptation with Tokens.}
TCA~\citep{wang2025tca} explores token condensation for test-time adaptation in vision-language models (CLIP). Our work differs in targeting standard ViTs and making token reduction itself robust, rather than using tokens for model adaptation.

\section{Method}
\label{sec:method}

\subsection{Background: Token Merging (ToMe)}

ToMe~\citep{bolya2023token} reduces the token count in each transformer layer by merging similar tokens using bipartite soft matching on key cosine similarity. Given a fixed merging ratio $r$ (tokens merged per layer), ToMe partitions tokens into two sets via alternating assignment, computes pairwise key similarities between sets, and greedily merges the $r$ most similar pairs by averaging their features (weighted by token size). This is applied identically at every layer regardless of input difficulty.

\subsection{Entropy-Guided Token Merging (\ours{})}

\ours{} modifies ToMe with a single core component: an \textbf{adaptive merging ratio} driven by attention entropy. The key insight is that CLS-row attention entropy increases under distribution shift, providing a lightweight signal to reduce merging when merging decisions become unreliable.

\paragraph{Step 1: Calibration.}
Given a small set of clean images ($n{=}500$), we run forward passes through the ViT with standard ToMe and record per-layer CLS-row attention entropy statistics. At each transformer layer $l$, we compute:
\begin{equation}
H_l = -\frac{1}{N_h} \sum_{h=1}^{N_h} \sum_{j} A^{(h)}_{\text{CLS},j} \log A^{(h)}_{\text{CLS},j}
\end{equation}
where $A^{(h)}_{\text{CLS},j}$ is the attention weight from the CLS token to token $j$ in head $h$, and $N_h$ is the number of attention heads. This requires only the CLS row of the attention matrix ($1 \times N$ per head), which is already computed during the forward pass. We store the per-layer mean $\bar{H}_l$ and standard deviation $\sigma_l$ from the calibration set.

\paragraph{Step 2: Inference-Time Adaptation.}
At inference, for each input and each layer, we compute $H_l$ and derive an entropy deviation score:
\begin{equation}
\delta_l = \max\!\left(0,\; \frac{H_l - \bar{H}_l}{\sigma_l}\right)
\end{equation}
The adaptive merging ratio is:
\begin{equation}
r_l = \left\lfloor r_0 \cdot \max\!\left(\alpha,\; \exp(-\beta \cdot \delta_l)\right) \right\rceil
\label{eq:adaptive_ratio}
\end{equation}
where $r_0$ is the base merging ratio (e.g., $r_0{=}8$ for DeiT-S, $r_0{=}16$ for DeiT-B), $\alpha \in (0,1)$ is a minimum ratio floor ensuring some merging always occurs, $\beta > 0$ controls sensitivity to entropy deviation, and $\lfloor \cdot \rceil$ denotes rounding to the nearest integer. When entropy is normal ($\delta_l \approx 0$), the ratio is close to $r_0$; when entropy is elevated (shift detected), the ratio decreases exponentially.

\paragraph{Computational Overhead.}
The CLS-row entropy computation adds $O(N)$ per head per layer, where $N$ is the current token count. Since attention scores are already computed in the forward pass, the additional cost is one entropy computation over a pre-existing softmax output. In our experiments, this adds $<2\%$ latency overhead when computed at every third layer (stride-3 variant), and approximately 35\% when computed at every layer due to the overhead of attention hook registration. We discuss this tradeoff in Section~\ref{sec:experiments}.

\paragraph{Design Rationale.}
The CLS token aggregates global information and its attention distribution reflects overall input uncertainty. Under corruption, CLS attention becomes more diffuse as the model struggles to identify reliable tokens~\citep{guo2023robustifying}. By monitoring this entropy, \ours{} naturally allocates more computation (less merging) to harder inputs. This is a form of input-adaptive computation that emerges from the shift-detection signal without explicit difficulty estimation.

\section{Experiments}
\label{sec:experiments}

\subsection{Setup}

\paragraph{Models.}
We use pretrained DeiT-Small (22M parameters, 12 layers)~\citep{touvron2021training} and DeiT-Base (86M parameters, 12 layers) from the \texttt{timm} library. No fine-tuning is performed.

\paragraph{Datasets.}
We evaluate on a stratified 5,000-image subset of ImageNet-1K validation (5 images per class) for clean accuracy, and apply 12 of the 15 standard ImageNet-C corruptions~\citep{hendrycks2019benchmarking} at severity 5. The three omitted corruptions (glass blur, zoom blur, fog) were excluded due to computational budget constraints. All corruptions are generated on-the-fly using the \texttt{imagecorruptions} library. We acknowledge that evaluation on the full 50K validation set and all 15 corruptions would strengthen confidence in generalizability (see Section~\ref{sec:discussion}).

\paragraph{Token Reduction Methods.}
We compare five methods, all training-free:
\begin{itemize}
    \item \textbf{No Reduction}: Full ViT inference (accuracy upper bound).
    \item \textbf{ToMe}~\citep{bolya2023token}: Fixed ratio $r{=}8$ (DeiT-S) / $r{=}16$ (DeiT-B).
    \item \textbf{Random Drop}: Randomly removes $r$ tokens per layer (CLS preserved). Lower bound on intelligent reduction.
    \item \textbf{EViT-style} (training-free): Applies EViT's token selection criterion~\citep{liang2022evit}---keep top-$k$ tokens by CLS attention, fuse rest---\emph{without} the fine-tuning step. Performance gap versus fine-tuned EViT is expected and discussed.
    \item \textbf{\ours{}} (ours): Entropy-guided adaptive merging with $\alpha{=}0.3$, $\beta{=}1.0$, applied at all 12 layers. Run with 3 calibration seeds (42, 43, 44); we report mean $\pm$ std.
\end{itemize}

\paragraph{Metrics.}
\begin{itemize}
    \item \textbf{Clean Accuracy}: Top-1 on clean ImageNet validation subset.
    \item \textbf{Mean Corrupt Accuracy}: Average top-1 across 12 corruption types at severity 5.
    \item \textbf{Relative Accuracy Drop}: $\text{RelDrop} = (\text{Acc}_{\text{clean}} - \text{Acc}_{\text{corrupt}}) / \text{Acc}_{\text{clean}}$, computed per-method (each method's own clean accuracy as reference).
    \item \textbf{Degradation Gap}: $\text{Gap} = \text{RelDrop}_{\text{method}} - \text{RelDrop}_{\text{no-reduction}}$. Positive values indicate the method degrades \emph{more} than the uncompressed baseline under corruption. This is the core quantity measuring disproportionate degradation from token reduction.
    \item \textbf{Throughput}: Images/second on a single NVIDIA RTX A6000 GPU.
\end{itemize}

\paragraph{Implementation.}
All experiments use batch size 128 (DeiT-S) or 64 (DeiT-B), 4 data workers, and PyTorch 2.1 with CUDA. Throughput is measured with 20 warmup + 100 timed batches using \texttt{torch.cuda.synchronize()}.

\subsection{Main Results}

\paragraph{DeiT-S ($r{=}8$): Conservative Merging.}
Table~\ref{tab:main_deit_s} presents results on DeiT-S. The disproportionate degradation gap from ToMe is \textbf{1.47\%}---below our pre-registered hypothesis threshold of 2\%. This indicates that with conservative merging ($r{=}8$, merging 8 of ${\sim}197$ tokens per layer), token merging does not substantially harm robustness on DeiT-S. \ours{} reduces the gap to 0.29\% (80.5\% recovery) but at the cost of throughput: \ours{}'s throughput (1,221 img/s) is actually \emph{lower} than the uncompressed baseline (1,444 img/s), because the entropy computation overhead exceeds the savings from merging ${\sim}6$ tokens per layer on average.

\begin{table}[t]
\caption{Results on DeiT-Small with $r{=}8$. ``Gap'' measures disproportionate degradation relative to No Reduction (Gap $= \text{RelDrop}_{\text{method}} - \text{RelDrop}_{\text{no-red.}}$). $\uparrow$~higher is better, $\downarrow$~closer to 0 is better. \ours{} reports mean $\pm$ std over 3 calibration seeds.}
\label{tab:main_deit_s}
\centering
\small
\begin{tabular}{lccccr}
\toprule
Method & Clean $\uparrow$ & Corrupt $\uparrow$ & Rel.\ Drop & Gap $\downarrow$ & Thru.\ $\uparrow$ \\
\midrule
No Reduction    & 77.86 & 43.14 & 44.59\% & ---   & 1,444 \\
Random Drop     & 76.12 & 39.24 & 48.44\% & +3.85 & 1,981 \\
EViT-style$^\dag$      & 76.80 & 42.21 & 45.04\% & +0.45 & 1,801 \\
ToMe $r{=}8$   & 77.16 & 41.62 & 46.06\% & +1.47 & 1,573 \\
\textbf{\ours{}} & 77.64{\scriptsize$\pm$0.00} & 42.80{\scriptsize$\pm$0.07} & 44.88\% & \textbf{+0.29} & 1,221 \\
\bottomrule
\end{tabular}
\vspace{-2mm}
\begin{flushleft}
{\footnotesize $^\dag$Training-free variant of EViT; gap versus fine-tuned EViT is expected.\\
Accuracies reported as percentages (\%). Throughput in images/second.}
\end{flushleft}
\end{table}

\paragraph{DeiT-B ($r{=}16$): Aggressive Merging.}
Table~\ref{tab:main_deit_b} shows that DeiT-B with aggressive merging ($r{=}16$) reveals a dramatically different picture. ToMe's degradation gap is \textbf{11.57\%}---nearly 8$\times$ larger than DeiT-S and well above the 2\% threshold. Random dropping is even worse (22.34\% gap), confirming that intelligent merging partially mitigates degradation but is insufficient. EViT-style shows a moderate gap (2.65\%).

\ours{} \textbf{fully recovers} the disproportionate degradation, achieving a gap of $-0.20\%$ (slightly \emph{better} than no reduction in terms of relative drop). This comes at a throughput cost: \ours{} achieves 510 img/s versus ToMe's 773 img/s. However, \ours{} still provides a $1.15\times$ speedup over the uncompressed baseline (442 img/s), demonstrating that adaptive merging can simultaneously provide both efficiency and robustness benefits on larger models.

\begin{table}[t]
\caption{Results on DeiT-Base with $r{=}16$. Aggressive merging causes severe disproportionate degradation (ToMe gap: +11.57\%), which \ours{} fully recovers (gap: $-0.20\%$).}
\label{tab:main_deit_b}
\centering
\small
\begin{tabular}{lccccr}
\toprule
Method & Clean $\uparrow$ & Corrupt $\uparrow$ & Rel.\ Drop & Gap $\downarrow$ & Thru.\ $\uparrow$ \\
\midrule
No Reduction    & 79.30 & 50.62 & 36.17\% & ---    & 442 \\
Random Drop     & 62.50 & 25.93 & 58.51\% & +22.34 & 936 \\
EViT-style$^\dag$      & 78.66 & 48.12 & 38.82\% & +2.65  & 555 \\
ToMe $r{=}16$  & 75.00 & 39.19 & 47.74\% & +11.57 & 773 \\
\textbf{\ours{}} & 76.77{\scriptsize$\pm$0.02} & \textbf{49.16}{\scriptsize$\pm$0.03} & 35.97\% & \textbf{$-$0.20} & 510 \\
\bottomrule
\end{tabular}
\end{table}

\paragraph{Scaling Behavior.}
The results reveal a clear scaling pattern: the disproportionate degradation from token merging grows with model size and merging aggressiveness. With DeiT-S ($r{=}8$), only 4\% of tokens are merged per layer and the gap is 1.47\%. With DeiT-B ($r{=}16$), approximately 8\% of tokens are merged per layer and the gap jumps to 11.57\%. This suggests that \ours{}'s adaptive mechanism becomes increasingly valuable as models grow larger and merging ratios more aggressive---precisely the regime where token merging provides the largest efficiency gains.

\subsection{Per-Corruption Analysis}

Table~\ref{tab:per_corruption_b} breaks down results by corruption type for DeiT-B, where the effects are most pronounced.

\begin{table}[t]
\caption{Per-corruption degradation gap on DeiT-B ($r{=}16$). Gap $=$ RelDrop$_{\text{method}} -$ RelDrop$_{\text{no-red.}}$. Negative \ours{} gap (bold) means it degrades \emph{less} than no reduction.}
\label{tab:per_corruption_b}
\centering
\small
\begin{tabular}{lcccc}
\toprule
Corruption & No Red.\ Acc & ToMe Acc & \ours{} Acc & ToMe Gap / \ours{} Gap \\
\midrule
\multicolumn{5}{l}{\textit{Noise}} \\
Gaussian noise      & 48.50 & 31.88 & 45.96 & +18.65 / +1.30 \\
Shot noise          & 48.38 & 32.16 & 45.73 & +18.13 / +1.44 \\
Impulse noise       & 48.70 & 32.26 & 46.35 & +18.40 / +1.04 \\
\midrule
\multicolumn{5}{l}{\textit{Blur}} \\
Defocus blur        & 30.92 & 22.60 & 29.16 & +8.86 / +1.01 \\
Motion blur         & 36.82 & 26.74 & 35.73 & +10.78 / \textbf{$-$0.10} \\
\midrule
\multicolumn{5}{l}{\textit{Weather}} \\
Snow                & 58.78 & 47.80 & 58.59 & +10.39 / \textbf{$-$2.20} \\
Frost               & 62.36 & 47.42 & 61.53 & +15.41 / \textbf{$-$1.51} \\
Brightness          & 72.04 & 64.62 & 69.55 & +4.68 / +0.25 \\
\midrule
\multicolumn{5}{l}{\textit{Digital}} \\
Contrast            & 73.46 & 66.22 & 71.51 & +4.34 / \textbf{$-$0.51} \\
Elastic transform   & 26.00 & 17.96 & 26.05 & +8.84 / \textbf{$-$1.14} \\
Pixelate            & 43.70 & 30.80 & 42.62 & +14.04 / \textbf{$-$0.41} \\
JPEG compression    & 57.74 & 49.86 & 57.11 & +6.33 / \textbf{$-$1.57} \\
\bottomrule
\end{tabular}
\end{table}

ToMe shows substantial degradation across all corruption types, with noise corruptions causing the largest gaps (18--19\%). \ours{} reduces the gap to near zero or negative across all corruptions. In 7 of 12 corruptions, \ours{}'s gap is negative, meaning it actually preserves robustness \emph{better} than the uncompressed model---likely because the reduced merging under corruption preserves discriminative tokens that would otherwise be merged.

\begin{figure}[t]
\centering
\includegraphics[width=0.95\linewidth]{figures/figure3_per_corruption.pdf}
\caption{EGTM recovery over ToMe per corruption type on DeiT-S (severity 5). Bars show the percentage of ToMe's disproportionate degradation gap that \ours{} recovers. Colors indicate entropy direction: blue = entropy increases under corruption, orange = entropy decreases, gray = neutral. Dashed line marks 50\% recovery. \ours{} achieves $>$60\% recovery on most corruption types.}
\label{fig:per_corruption}
\end{figure}

\subsection{Entropy-Severity Correlation}

We analyze the relationship between CLS-row attention entropy and corruption severity using Spearman rank correlation on DeiT-S across 5 severity levels (Table~\ref{tab:entropy_corr}). Of 8 tested corruption types, 5 show statistically significant positive correlations ($p < 0.05$): Gaussian noise ($\rho{=}1.0$), shot noise ($\rho{=}1.0$), impulse noise ($\rho{=}0.9$), snow ($\rho{=}1.0$), and brightness ($\rho{=}0.9$). Three corruptions---defocus blur, motion blur, and frost---show positive but \emph{non-significant} correlations ($\rho{=}0.7$, $p{=}0.188$). The non-significance stems from the small sample size (5 severity levels) and non-monotonic entropy behavior at intermediate severities. Importantly, all 8 tested corruptions show positive $\rho$, suggesting the general trend holds even where statistical significance is lacking.

\paragraph{Note on Success Criterion.} Our proposal set a success criterion of $\rho > 0.5$ for $>$10 of 15 corruption types. Only 8 of 12 evaluated corruptions were tested for entropy-severity correlation (the remaining 4 were omitted due to compute constraints), and 3 of the original 15 corruption types were excluded entirely from evaluation. Therefore, this criterion cannot be fully assessed. Among the 8 tested, all show $\rho \geq 0.7$ (all exceed the 0.5 threshold), and 5 of 8 reach statistical significance. While these results are encouraging, they do not constitute a definitive test of the proposed criterion. Full evaluation across all 15 corruption types with finer severity granularity is needed.

\begin{table}[t]
\caption{Entropy-severity correlation (Spearman $\rho$) on DeiT-S across 5 severity levels. $^*$: $p < 0.05$; $^\dagger$: $p \geq 0.05$ (not statistically significant at the 0.05 level).}
\label{tab:entropy_corr}
\centering
\small
\begin{tabular}{lcc}
\toprule
Corruption & Spearman $\rho$ & $p$-value \\
\midrule
Gaussian noise$^*$   & 1.00 & $< 0.001$ \\
Shot noise$^*$       & 1.00 & $< 0.001$ \\
Impulse noise$^*$    & 0.90 & 0.037 \\
Snow$^*$             & 1.00 & $< 0.001$ \\
Brightness$^*$       & 0.90 & 0.037 \\
Defocus blur$^\dagger$ & 0.70 & 0.188 \\
Motion blur$^\dagger$  & 0.70 & 0.188 \\
Frost$^\dagger$        & 0.70 & 0.188 \\
\bottomrule
\end{tabular}
\end{table}

\paragraph{Entropy Direction Analysis.}
Among the 12 corruptions evaluated for accuracy, the per-corruption analysis shows that most corruptions either increase or maintain CLS-row entropy relative to clean inputs. Brightness shows a decrease in mean entropy ($\Delta = -1.25\sigma$), classified as a ``decrease'' direction---a predicted failure mode where the model becomes overconfident rather than uncertain. Despite this, \ours{} still performs well on brightness (gap: $-0.22\%$ on DeiT-S) because the entropy floor parameter $\alpha$ ensures a minimum merging reduction is always applied.

\subsection{Adaptive Merging Behavior}

Figure~\ref{fig:adaptive_ratio} shows how \ours{} adapts its merging ratio. On DeiT-S, the average merging ratio drops from 6.2 (clean) to 5.5 (corrupted)---an 11\% reduction. On DeiT-B, the drop is more pronounced: from 14.3 (clean) to 11.7 (corrupted)---an 18\% reduction. The adaptation is corruption-dependent: elastic transform triggers the strongest reduction (avg $r \approx 9.8$ on DeiT-B vs.\ base $r_0{=}16$), while brightness triggers minimal adaptation ($r \approx 14.4$), consistent with the entropy direction analysis.

\begin{figure}[t]
\centering
\includegraphics[width=0.8\linewidth]{figures/figure2_adaptive_ratio.pdf}
\caption{Average merging ratio for \ours{} on clean vs.\ corrupted inputs across layers (DeiT-S). The method automatically merges fewer tokens under corruption.}
\label{fig:adaptive_ratio}
\end{figure}

\subsection{Ablation Study}
\label{sec:ablation}

Table~\ref{tab:ablation} presents ablation results on DeiT-S with three representative corruptions (Gaussian noise, motion blur, snow) at severity 5.

\begin{table}[t]
\caption{Ablation study on DeiT-S. ``Corrupt Acc'' is the mean over Gaussian noise, motion blur, and snow at severity 5. Avg.\ $r$ shows mean tokens merged per layer.}
\label{tab:ablation}
\centering
\small
\begin{tabular}{lcccr}
\toprule
Variant & Clean Acc & Corrupt Acc & Avg.\ $r$ & Thru.\ (img/s) \\
\midrule
ToMe $r{=}8$ (base)         & 77.16 & 39.21 & 8.0 & 1,573 \\
ToMe $r{=}4$ (conservative) & 77.72 & 40.19 & 4.0 & 1,340 \\
Importance prot.\ only       & 77.82 & 40.06 & 8.0 &   718 \\
\ours{} (CLS entropy)       & 77.64 & 40.15 & 5.4 & 1,221 \\
\ours{} + importance prot.  & 77.64 & 40.15 & 5.4 & 1,219 \\
\ours{} (full entropy)      & 77.64 & 39.61 & 5.0 & 1,119 \\
\ours{} (CLS variance)      & 77.54 & 40.27 & 3.8 &   975 \\
\bottomrule
\end{tabular}
\end{table}

Key findings:
\begin{enumerate}
    \item \textbf{Fixed conservative ratio is competitive on DeiT-S.} ToMe $r{=}4$ achieves 40.19\% corrupt accuracy with 1,340 img/s throughput, compared to \ours{}'s 40.15\% at 1,221 img/s. On DeiT-S with small merging ratios, simply halving the ratio is a simpler and faster alternative.
    \item \textbf{Adaptive ratio is the primary contributor.} Adding importance protection to \ours{} provides negligible additional benefit (40.15\% in both cases), suggesting the adaptive ratio alone captures the key mechanism.
    \item \textbf{CLS-row entropy is the best proxy.} The CLS-row entropy variant achieves the best accuracy-throughput tradeoff. Full attention matrix entropy merges too aggressively (avg $r{=}5.0$) with lower throughput. CLS variance is overly conservative (avg $r{=}3.8$) and slowest.
    \item \textbf{The value of adaptiveness emerges on larger models.} The ablation on DeiT-S shows marginal differences because the degradation gap is small (1.47\%). On DeiT-B (Table~\ref{tab:main_deit_b}), where the gap is 11.57\%, \ours{}'s adaptive behavior is critical---a fixed conservative ratio cannot match \ours{}'s ability to merge aggressively on clean inputs while preserving tokens under corruption.
    \item \textbf{Stride-based entropy computation improves DeiT-S throughput.} Computing entropy at every 4th layer (stride-4 variant) instead of all 12 layers improves DeiT-S throughput from 1,221 to 1,372 img/s---closer to the uncompressed baseline (1,444 img/s)---with comparable robustness (mean corrupt accuracy: 42.2\% stride-4 vs.\ 42.8\% all-layers). For DeiT-S, the stride-4 variant offers a better practical tradeoff.
\end{enumerate}

\subsection{Hyperparameter Sensitivity}

Figure~\ref{fig:sensitivity} shows sensitivity to $\alpha$ (minimum ratio floor) and $\beta$ (entropy sensitivity) on DeiT-S.

\begin{figure}[t]
\centering
\includegraphics[width=0.95\linewidth]{figures/figure5_sensitivity.pdf}
\caption{Hyperparameter sensitivity on DeiT-S. Both clean and corrupt accuracy are robust to $\alpha \in [0.1, 0.5]$ and $\beta \in [0.3, 3.0]$.}
\label{fig:sensitivity}
\end{figure}

\ours{} is robust to hyperparameter choices: corrupt accuracy varies by only 0.8 percentage points across the full $\alpha$ range $[0.1, 0.7]$ and 0.02 percentage points across $\beta \in [0.3, 3.0]$. The default $\alpha{=}0.3$, $\beta{=}1.0$ performs well across both settings. This stability arises because the $z$-score normalization already calibrates the entropy signal, making the exponential decay relatively insensitive to the exact $\beta$ value.

\subsection{Entropy Distribution Visualization}

Figure~\ref{fig:entropy_dist} shows the distribution of CLS-row attention entropy across layers for clean versus Gaussian noise-corrupted inputs on DeiT-S. Corrupted inputs consistently show higher entropy in later layers (6--11), where the model makes higher-level feature decisions. This layer-dependent pattern explains why \ours{}'s per-layer adaptation is effective: it preserves tokens precisely in the layers where corruption disrupts attention the most.

\begin{figure}[t]
\centering
\includegraphics[width=0.95\linewidth]{figures/figure1_entropy_distribution.pdf}
\caption{CLS-row attention entropy distribution across DeiT-S layers for clean (blue) vs.\ Gaussian noise severity 5 (red). Corrupted inputs show elevated entropy, particularly in later layers.}
\label{fig:entropy_dist}
\end{figure}

\section{Discussion and Limitations}
\label{sec:discussion}

\paragraph{When \ours{} Helps and When It Does Not.}
Our results reveal a clear dichotomy. On DeiT-B with aggressive merging ($r{=}16$), token merging causes severe disproportionate degradation (11.57\% gap) and \ours{} provides substantial value by fully recovering it. On DeiT-S with conservative merging ($r{=}8$), the degradation gap is modest (1.47\%, below our 2\% hypothesis threshold), \ours{}'s throughput overhead exceeds its merging savings, and a simple fixed $r{=}4$ achieves comparable robustness. \textbf{We recommend \ours{} for larger models with aggressive merging ratios}, where the efficiency-robustness tradeoff is most severe and adaptive behavior provides clear benefits.

\paragraph{Throughput Tradeoffs.}
\ours{}'s entropy computation adds overhead that is proportionally more costly for smaller, faster models. On DeiT-S, \ours{} (1,221 img/s) is 15\% slower than the uncompressed baseline (1,444 img/s), providing negative net efficiency. On DeiT-B, \ours{} achieves 510 img/s---a $1.15\times$ speedup over the baseline (442 img/s), modest compared to ToMe's $1.75\times$ but with dramatically better robustness. The throughput gain fraction (relative to ToMe's gain over baseline) is $(510-442)/(773-442) = 20.5\%$. Future work could reduce overhead via strided entropy computation (every $k$-th layer), fused CUDA kernels, or amortized calibration.

\paragraph{SATA Comparison.}
Our re-implementation of SATA~\citep{nikzad2025sata} using simplified spatial grouping yielded substantially degraded clean accuracy (46.7\% on DeiT-S vs.\ the 77.9\% expected for an unmodified model), indicating our spatial grouping approximation was insufficiently faithful to the original method. We exclude SATA from main comparisons to avoid unfair evaluation. A proper comparison using the authors' official implementation is important future work.

\paragraph{Evaluation Scope.}
Our evaluation uses a 5K-image subset (not the full 50K ImageNet validation), 12 of 15 corruption types, severity 5 only (not 3+5), and does not include ImageNet-R or ImageNet-Sketch. While the trends are clear and consistent across corruptions, confirming these results on the complete evaluation suite is necessary before drawing deployment recommendations. The evaluation on only severity 5 means we cannot characterize how the degradation gap scales with corruption severity.

\paragraph{Failure Modes.}
We identify two failure modes: (1) corruptions that \emph{decrease} attention entropy (e.g., brightness), where the model becomes overconfident and \ours{} does not trigger adaptation; and (2) calibration domain mismatch, where deployment on a domain substantially different from calibration data (e.g., medical imaging) could miscalibrate the entropy statistics. The $\alpha$ floor parameter and $z$-score normalization provide partial mitigation for both, but neither is a complete solution. Adversarial perturbations designed to reduce attention entropy could also potentially circumvent \ours{}'s mechanism.

\paragraph{Statistical Significance of Entropy Correlations.}
Three of eight tested entropy-severity correlations have $p = 0.188$, falling short of statistical significance. This is partly a sample-size limitation (only 5 severity levels per corruption), but also reflects genuine non-monotonicity in how some corruptions affect attention entropy at intermediate severities. We recommend that future work test on finer severity granularity to improve statistical power.

\section{Conclusion}

We presented the first systematic study of token reduction robustness under distribution shift, finding that aggressive token merging causes substantial disproportionate accuracy degradation---up to 11.6 percentage points additional relative accuracy drop on DeiT-B with $r{=}16$. We proposed \ours{}, a training-free method that adapts the merging ratio based on CLS-row attention entropy. On DeiT-B, \ours{} fully recovers the degradation gap while maintaining a $1.15\times$ speedup. We clearly characterized \ours{}'s limitations: on DeiT-S with conservative merging ($r{=}8$), the degradation is modest and \ours{}'s overhead exceeds its benefits.

Our findings suggest practical implications for deploying efficient ViTs, though we emphasize that these recommendations are \emph{provisional} pending validation on the full 50K ImageNet validation set, all 15 corruption types, multiple severity levels, and additional distribution shift benchmarks (ImageNet-R, ImageNet-Sketch). Subject to this caveat: practitioners using aggressive token merging on large models should be aware of potential robustness degradation and consider adaptive approaches like \ours{}. For small models with conservative merging, standard ToMe with a moderately reduced ratio may suffice. Future work should evaluate on the complete evaluation suite, explore combinations with complementary robustness methods like SATA, and develop lower-overhead shift-detection signals for efficient models.

\begin{thebibliography}{19}

\bibitem[Bolya et~al.(2023)]{bolya2023token}
Daniel Bolya, Cheng-Yang Fu, Xiaoliang Dai, Peizhao Zhang, Christoph Feichtenhofer, and Judy Hoffman.
\newblock Token merging: Your {ViT} but faster.
\newblock In \emph{International Conference on Learning Representations (ICLR)}, 2023.

\bibitem[Dosovitskiy et~al.(2021)]{dosovitskiy2021image}
Alexey Dosovitskiy, Lucas Beyer, Alexander Kolesnikov, Dirk Weissenborn, Xiaohua Zhai, Thomas Unterthiner, Mostafa Dehghani, Matthias Minderer, Georg Heigold, Sylvain Gelly, Jakob Uszkoreit, and Neil Houlsby.
\newblock An image is worth 16x16 words: Transformers for image recognition at scale.
\newblock In \emph{International Conference on Learning Representations (ICLR)}, 2021.

\bibitem[Guo et~al.(2023)]{guo2023robustifying}
Yong Guo, David Stutz, and Bernt Schiele.
\newblock Robustifying token attention for vision transformers.
\newblock In \emph{Proceedings of the IEEE/CVF International Conference on Computer Vision (ICCV)}, 2023.

\bibitem[Hendrycks \& Dietterich(2019)]{hendrycks2019benchmarking}
Dan Hendrycks and Thomas Dietterich.
\newblock Benchmarking neural network robustness to common corruptions and perturbations.
\newblock In \emph{International Conference on Learning Representations (ICLR)}, 2019.

\bibitem[Kim et~al.(2024)]{kim2024token}
Minchul Kim, Shangqian Gao, Yen-Chang Hsu, Yilin Shen, and Hongxia Jin.
\newblock Token fusion: Bridging the gap between token pruning and token merging.
\newblock In \emph{IEEE/CVF Winter Conference on Applications of Computer Vision (WACV)}, 2024.

\bibitem[Kong et~al.(2025)]{kong2025token}
Zhenglun Kong, Yize Li, Fanhu Zeng, Lei Xin, Shvat Messica, Xue Lin, Pu Zhao, Manolis Kellis, Hao Tang, and Marinka Zitnik.
\newblock Token reduction should go beyond efficiency in generative models---from vision, language to multimodality.
\newblock \emph{arXiv preprint arXiv:2505.18227}, 2025.

\bibitem[Liang et~al.(2022)]{liang2022evit}
Youwei Liang, Chongjian Ge, Zhan Tong, Yibing Song, Jue Wang, and Pengtao Xie.
\newblock Not all patches are what you need: Expediting vision transformers via token reorganizations.
\newblock In \emph{International Conference on Learning Representations (ICLR)}, 2022.

\bibitem[Lin et~al.(2026)]{lin2026entropy}
Shangqun Lin, Peng Lyu, Dongliang Liu, et~al.
\newblock Entropy-guided condensing for vision transformer.
\newblock \emph{International Journal of Computer Vision}, 134:86, 2026.

\bibitem[Nikzad et~al.(2025)]{nikzad2025sata}
Nick Nikzad, Yi Liao, Yongsheng Gao, and Jun Zhou.
\newblock {SATA}: Spatial autocorrelation token analysis for enhancing the robustness of vision transformers.
\newblock In \emph{Proceedings of the IEEE/CVF Conference on Computer Vision and Pattern Recognition (CVPR)}, 2025.

\bibitem[Rao et~al.(2021)]{rao2021dynamicvit}
Yongming Rao, Wenliang Zhao, Benlin Liu, Jiwen Lu, Jie Zhou, and Cho-Jui Hsieh.
\newblock {DynamicViT}: Efficient vision transformers with dynamic token sparsification.
\newblock In \emph{Advances in Neural Information Processing Systems (NeurIPS)}, 2021.

\bibitem[Shin et~al.(2025)]{shin2025atom}
J.~Shin, M.~Kang, Y.~Han, J.~Park, and L.-S. Kim.
\newblock {AToM}: Adaptive token merging for efficient acceleration of vision transformer.
\newblock \emph{IEEE Transactions on Computers}, 2025.

\bibitem[Touvron et~al.(2021)]{touvron2021training}
Hugo Touvron, Matthieu Cord, Matthijs Douze, Francisco Massa, Alexandre Sablayrolles, and Herv{\'e} J{\'e}gou.
\newblock Training data-efficient image transformers \& distillation through attention.
\newblock In \emph{International Conference on Machine Learning (ICML)}, 2021.

\bibitem[Wang et~al.(2025)]{wang2025tca}
Zixin Wang, Dong Gong, Sen Wang, Zi Huang, and Yadan Luo.
\newblock Is less more? {E}xploring token condensation as training-free test-time adaptation.
\newblock In \emph{IEEE/CVF International Conference on Computer Vision (ICCV)}, 2025.

\bibitem[Yin et~al.(2022)]{yin2022avit}
Hongxu Yin, Arash Vahdat, Jose~M. Alvarez, Arun Mallya, Jan Kautz, and Pavlo Molchanov.
\newblock {A-ViT}: Adaptive tokens for efficient vision transformer.
\newblock In \emph{IEEE/CVF Conference on Computer Vision and Pattern Recognition (CVPR)}, 2022.

\bibitem[Yoo et~al.(2024)]{yoo2024adaptiv}
Seungjae Yoo, Hangyeol Kim, and Joo-Young Kim.
\newblock {AdapTiV}: Sign-similarity based image-adaptive token merging for vision transformer acceleration.
\newblock In \emph{Proceedings of the 57th IEEE/ACM International Symposium on Microarchitecture (MICRO)}, 2024.

\bibitem[Zhai et~al.(2023)]{zhai2023stabilizing}
Shuangfei Zhai, Tatiana Likhomanenko, Etai Littwin, Dan Busbridge, Jason Ramapuram, Yizhe Zhang, Jiatao Gu, and Joshua~M. Susskind.
\newblock Stabilizing transformer training by preventing attention entropy collapse.
\newblock In \emph{International Conference on Machine Learning (ICML)}, 2023.

\bibitem[Zhang et~al.(2022)]{zhang2022delving}
Chongzhi Zhang, Mingyuan Zhang, Shanghang Zhang, Daisheng Jin, Qiang Zhou, Zhongang Cai, Haiyu Zhao, Xianglong Liu, and Ziwei Liu.
\newblock Delving deep into the generalization of vision transformers under distribution shifts.
\newblock In \emph{IEEE/CVF Conference on Computer Vision and Pattern Recognition (CVPR)}, 2022.

\end{thebibliography}

\end{document}
